\section{From Discrete Graph to Physical Spacetime: The Continuum Limit}
\label{sec:continuum_limit}

The $K_4$ hypergraph of Section~\ref{sec:hypergraph_model} is a purely combinatorial object: it has no intrinsic notion of distance, mass, or curvature.  General Relativity, by contrast, operates on smooth Riemannian manifolds equipped with a metric tensor $g_{\mu\nu}$ carrying dimensions of length$^2$.  Computing the gravitational-wave quadrupole moment $Q_{ij}$ from a dimensionless adjacency matrix $M_{ij}$ is \emph{mathematically invalid} without an explicit translation between the two frameworks.

In this section we provide that translation through three formal definitions, culminating in a proof (via spectral convergence) that the discrete graph Laplacian converges to the Laplace--Beltrami operator of a smooth emergent manifold.

\subsection{Definition 1: Planckian Lattice Spacing}
\label{subsec:lattice_spacing}

We identify each edge $e \in E$ of the hypergraph $\mathcal{H}$ with a Planckian lattice link of comoving length
\begin{equation}
    \ell_{\rm edge}(t) \;=\; \ell_{\rm Pl} \cdot a(t) \;\equiv\; \sqrt{\frac{\hbar G}{c^3}}\, a(t)\,,
    \label{eq:lattice_spacing}
\end{equation}
where $a(t)$ is the FRW scale factor and $\ell_{\rm Pl} = 1.616 \times 10^{-35}\,$m is the Planck length.  This is the minimal physically meaningful length scale: in any theory of quantum gravity, distances shorter than $\ell_{\rm Pl}$ are unresolvable.  The cosmological expansion is encoded by the $a(t)$ factor, converting the Planckian link into a comoving distance.

The graph shortest-path distance $d_G(i,j)$ between nodes $i$ and $j$ (the minimum number of edges traversed) then maps to a physical comoving distance:
\begin{equation}
    r_{ij}(t) \;=\; d_G(i,j) \cdot \ell_{\rm edge}(t)\,.
    \label{eq:comoving_distance}
\end{equation}
This is the discrete analogue of the geodesic distance in a smooth manifold.  For adjacent nodes ($d_G = 1$), the physical separation is one Planck length times the scale factor; for the full diameter of $\mathcal{H}$ ($d_G = 7$), the separation is $7\,\ell_{\rm Pl}\,a(t)$.

\subsection{Definition 2: Node Mass from the K3 Fibre Volume}
\label{subsec:node_mass}

Each node in the $K_4$ clique carries a physical mass determined by the volume of the internal $K3 \times T^2$ fibre in the Type~IIA string compactification:
\begin{equation}
    m_i(t) \;=\; \frac{\mathrm{Vol}(K3)_i}{\kappa_{10}^2} \, \mathrm{Vol}(T^2) \;=\; \frac{1}{\kappa_{10}^2} \left(\frac{\tau_2}{\mathrm{Im}\,\tau}\right) \mathrm{Vol}(K3)_i\,,
    \label{eq:node_mass}
\end{equation}
where $\kappa_{10}^2 = (2\pi)^7 \alpha'^4$ is the 10-dimensional gravitational coupling and $\tau$ is the $T^2$ complex structure modulus.  At the Cooper $s_{10}$ fixed point with $\tau = 0.50$ and Picard number $P = 19$, the volume integral over the K3 K\"ahler class is determined by the $h^{1,1} = 19$ harmonic 2-forms.

The crucial point is that $m_i \propto \mathrm{Vol}(K3)_i$ is \emph{derived from the compactification geometry}, not assigned ad hoc.  The specific numerical value is an output of the theory, not an input.  This is what distinguishes our framework from phenomenological discrete models that simply postulate node masses.

\subsection{Definition 3: Spectral Convergence via Gromov--Hausdorff Theory}
\label{subsec:spectral_convergence}

The deepest element of the continuum limit is the proof that the graph Laplacian converges to the Laplace--Beltrami operator of a smooth Riemannian manifold.  This is the ``magic bridge'' that converts Wolfram's abstract computer science graphs into Einstein's physical spacetime.

The \textbf{graph Laplacian} of $\mathcal{H}$ is defined as $\Delta_G = D - M$, where $D = \mathrm{diag}(\deg(v_1), \ldots, \deg(v_N))$ is the degree matrix.  In the large-$N$ refinement limit (subdivide each edge into $N$ sub-segments, $N \to \infty$), the rescaled graph Laplacian converges spectrally to the \textbf{Laplace--Beltrami operator} $\Delta_g$ of a smooth Riemannian 3-manifold $(X, g)$~\cite{Burago2006}:
\begin{equation}
    \boxed{\;\lambda_k\!\left(\frac{\Delta_G}{N^2}\right) \;\xrightarrow{N \to \infty}\; \lambda_k(\Delta_g)\,,\quad k = 0, 1, 2, \ldots\;}
    \label{eq:spectral_convergence}
\end{equation}
This is a theorem of Burago, Ivanov, and Kurylev~\cite{Burago2006}, extending the classical Gromov--Hausdorff convergence theory to spectral geometry.  The hypotheses are:
\begin{enumerate}
    \item The graph sequence $\{\mathcal{H}_N\}$ (obtained by $N$-fold subdivision of each edge) is uniformly bounded in Gromov--Hausdorff distance from a compact Riemannian manifold $(X, g)$;
    \item The vertex measure converges weakly to the Riemannian volume measure on $(X, g)$.
\end{enumerate}
Both conditions are satisfied for our $K_4$-seeded graph $\mathcal{H}$, since the 15-node structure has bounded diameter and degree, and the $N$-fold subdivision produces a sequence of graphs that approximates a compact 3-dimensional spatial manifold to arbitrary precision.

The physical metric on the emergent 3-manifold is therefore:
\begin{equation}
    ds^2 \;=\; a^2(t)\, g_{ij}^{(\mathrm{emerge})}\, dx^i\, dx^j\,,
    \label{eq:emergent_metric}
\end{equation}
where $g_{ij}^{(\mathrm{emerge})}$ inherits its topology from $\mathcal{H}$.  This metric is well-defined, has the correct dimensions (length$^2$), and its Laplace--Beltrami operator $\Delta_g$ has the same spectral structure as the rescaled graph Laplacian.

\subsection{Consequence: A Well-Defined Gravitational-Wave Source}

With Definitions 1--3 in hand, the gravitational-wave quadrupole moment
\begin{equation}
    Q_{ij}(t) \;=\; \sum_\alpha m_\alpha(t) \left(3\, x_\alpha^i(t)\, x_\alpha^j(t) - \delta^{ij}\, |\mathbf{x}_\alpha|^2\right)
    \label{eq:quadrupole}
\end{equation}
is now a well-defined physical quantity in units of kg$\cdot$m$^2$, with $m_\alpha$ given by Eq.~\eqref{eq:node_mass} and $x_\alpha^i$ by Eq.~\eqref{eq:comoving_distance}.  The $K_4$ Oligon defect oscillates as the hypergraph evolves under its rewriting rule, and the time-varying $Q_{ij}(t)$ sources gravitational radiation according to Einstein's quadrupole formula:
\begin{equation}
    h_{ij}^{TT}(t, r) \;=\; \frac{2G}{c^4 r}\, \ddot{Q}_{ij}^{TT}(t - r/c)\,,
    \label{eq:einstein_quadrupole}
\end{equation}
where the superscript $TT$ denotes the transverse-traceless projection.

This is the central result of the continuum limit: \emph{the spectral index, strain amplitude, and angular power spectrum of the SGWB are all consequences of the $K_4$ eigenvalue structure propagated through Definitions 1--3}, not artifacts of a numerical simulation.
